\newcommand{\figuraFlujoIncompresible}{%
\begin{figure}[htbp]
    \centering
    \begin{tikzpicture}[>=Latex,font=\small]
        \node[font=\bfseries] at (6.55,5.95)
            {Deformación de un elemento material incompresible};

        % Estado inicial.
        \coordinate (A) at (0.65,1.15);
        \coordinate (B) at (4.05,1.15);
        \coordinate (C) at (4.05,4.35);
        \coordinate (D) at (0.65,4.35);
        \coordinate (E) at (1.45,1.75);
        \coordinate (F) at (4.85,1.75);
        \coordinate (G) at (4.85,4.95);
        \coordinate (H) at (1.45,4.95);

        \fill[blue!5] (A)--(B)--(C)--(D)--cycle;
        \fill[blue!10] (B)--(F)--(G)--(C)--cycle;
        \fill[blue!16] (D)--(C)--(G)--(H)--cycle;
        \draw[very thick] (A)--(B)--(C)--(D)--cycle;
        \draw[very thick] (B)--(F)--(G)--(C);
        \draw[very thick] (C)--(G)--(H)--(D);
        \draw[very thick,dashed] (A)--(E)--(F);
        \draw[very thick,dashed] (E)--(H);

        \foreach \x/\y in {
            1.05/1.55,1.85/1.42,2.70/1.62,3.50/1.40,
            1.10/2.25,1.95/2.55,2.85/2.22,3.62/2.62,
            1.18/3.10,2.05/3.45,2.92/3.05,3.70/3.52,
            1.22/3.92,2.18/4.10,3.05/3.88,3.75/4.18,
            4.34/2.05,4.48/2.85,4.40/3.68,1.70/4.58,2.75/4.62}
        {\shade[ball color=blue!65] (\x,\y) circle (0.075);}

        \node[font=\bfseries] at (2.72,0.62) {estado inicial, $t$};
        \node[align=center] at (2.72,-0.05)
            {$M(t)=M_0$ \qquad $V(t)=V_0$};

        % Evolución temporal.
        \draw[very thick,green!45!black,-{Latex[length=4mm]}]
            (5.25,2.95)--node[above]{$t\rightarrow t+dt$}(7.15,2.95);

        % Estado deformado con el mismo volumen.
        \coordinate (Ap) at (7.65,1.15);
        \coordinate (Bp) at (11.35,1.15);
        \coordinate (Cp) at (12.15,4.35);
        \coordinate (Dp) at (8.45,4.35);
        \coordinate (Ep) at (8.45,1.75);
        \coordinate (Fp) at (12.15,1.75);
        \coordinate (Gp) at (12.95,4.95);
        \coordinate (Hp) at (9.25,4.95);

        \fill[blue!5] (Ap)--(Bp)--(Cp)--(Dp)--cycle;
        \fill[blue!10] (Bp)--(Fp)--(Gp)--(Cp)--cycle;
        \fill[blue!16] (Dp)--(Cp)--(Gp)--(Hp)--cycle;
        \draw[very thick] (Ap)--(Bp)--(Cp)--(Dp)--cycle;
        \draw[very thick] (Bp)--(Fp)--(Gp)--(Cp);
        \draw[very thick] (Cp)--(Gp)--(Hp)--(Dp);
        \draw[very thick,dashed] (Ap)--(Ep)--(Fp);
        \draw[very thick,dashed] (Ep)--(Hp);

        \foreach \x/\y in {
            8.05/1.52,8.88/1.42,9.75/1.58,10.62/1.40,
            8.30/2.22,9.18/2.48,10.05/2.22,10.88/2.58,
            8.55/3.02,9.45/3.38,10.32/3.05,11.20/3.48,
            8.78/3.88,9.72/4.08,10.62/3.85,11.48/4.15,
            11.72/2.02,11.92/2.82,12.10/3.65,9.45/4.58,10.55/4.62}
        {\shade[ball color=blue!65] (\x,\y) circle (0.075);}

        % Movimiento de cizalla: cambia la forma, no el volumen.
        \draw[very thick,green!45!black,-{Latex[length=3mm]}]
            (8.05,5.28)--(9.05,5.28);
        \draw[very thick,green!45!black,-{Latex[length=3mm]}]
            (11.35,5.28)--(12.35,5.28);
        \draw[very thick,green!45!black,-{Latex[length=3mm]}]
            (8.65,1.00)--(7.65,1.00);
        \draw[very thick,green!45!black,-{Latex[length=3mm]}]
            (11.85,1.00)--(10.85,1.00);
        \node[green!40!black] at (10.20,5.42)
            {cizalla};

        \node[font=\bfseries] at (10.28,0.62)
            {estado deformado, $t+dt$};
        \node[align=center] at (10.28,-0.05)
            {$M(t+dt)=M_0$ \qquad $V(t+dt)=V_0$};

        \node[align=center,text=black!75] at (6.55,-0.88)
            {$\displaystyle\rho=\frac{M_0}{V_0}=\rho_0$
             \qquad
             $\displaystyle\frac{D\rho}{Dt}=0$
             \qquad$\Longleftrightarrow$\qquad
             $\displaystyle\vec{\nabla}\cdot\vec{v}=0$};
    \end{tikzpicture}
    \caption{Evolución de un elemento material en un flujo incompresible. El
    elemento puede cambiar de forma por cizalla, pero conserva su masa, su
    volumen y, por tanto, su densidad. El campo de velocidad tiene divergencia
    nula.}
    \label{fig:flujo-incompresible}
\end{figure}%
}